\section{State Representation in Sequential Decision-Making}

\subsection{Overview}

In supervised and self-supervised learning, data is typically assumed to be static.  
In contrast, reinforcement learning (RL) models \emph{sequential interaction} between an agent and an environment.

\medskip

\noindent
\textbf{Key shift.}  
Data is no longer given; it is generated through actions.  
Learning must therefore account for feedback loops between decisions and future observations.

\subsection{Markov Decision Processes}

\begin{definition}[Markov Decision Process (MDP)]
A Markov Decision Process is a tuple
\[
(\mathcal{S}, \mathcal{A}, P, R, \gamma),
\]
where:
\begin{itemize}
\item $\mathcal{S}$ is a state space
\item $\mathcal{A}$ is an action space
\item $P(s' \mid s,a)$ is the transition probability
\item $R(s,a)$ is the reward function
\item $\gamma \in [0,1)$ is the discount factor
\end{itemize}
\end{definition}

\begin{definition}[Markov Property]
An environment satisfies the Markov property if
\[
P(s_{t+1} \mid s_t, a_t, s_{t-1}, a_{t-1}, \dots)
= P(s_{t+1} \mid s_t, a_t).
\]
\end{definition}

\subsection{Policies and Returns}

\begin{definition}[Policy]
A policy is a mapping
\[
\pi(a \mid s),
\]
giving the probability of taking action $a$ in state $s$.
\end{definition}

\begin{definition}[Return]
The return from time $t$ is
\[
G_t = \sum_{k=0}^{\infty} \gamma^k R(s_{t+k}, a_{t+k}).
\]
\end{definition}

\begin{definition}[Value Function]
The state-value function of a policy $\pi$ is
\[
V^\pi(s) = \mathbb{E}_\pi[G_t \mid s_t = s].
\]
\end{definition}

\begin{definition}[Action-Value Function]
The action-value function is
\[
Q^\pi(s,a) = \mathbb{E}_\pi[G_t \mid s_t = s, a_t = a].
\]
\end{definition}

\subsection{Bellman Equations}

\begin{proposition}[Bellman Expectation Equation]
For any policy $\pi$,
\[
V^\pi(s) = \sum_a \pi(a \mid s)
\left[
R(s,a) + \gamma \sum_{s'} P(s' \mid s,a) V^\pi(s')
\right].
\]
\end{proposition}

\begin{proposition}[Bellman Optimality Equation]
The optimal value function satisfies
\[
V^*(s) = \max_a
\left[
R(s,a) + \gamma \sum_{s'} P(s' \mid s,a) V^*(s')
\right].
\]
\end{proposition}

\begin{remark}
The Bellman equations express value functions as fixed points of nonlinear operators.
\end{remark}

\subsection{Bellman Fixed-Point Theorem}

\begin{theorem}[Contraction of the Bellman Operator]
Let $\mathcal{T}$ be the Bellman optimality operator:
\[
(\mathcal{T}V)(s)
= \max_a \left[ R(s,a) + \gamma \sum_{s'} P(s' \mid s,a) V(s') \right].
\]
Then for $\gamma \in [0,1)$, $\mathcal{T}$ is a contraction in the sup norm:
\[
\|\mathcal{T}V - \mathcal{T}W\|_\infty
\le \gamma \|V - W\|_\infty.
\]
\end{theorem}

\begin{proof}
For any $s$,
\[
|(\mathcal{T}V)(s) - (\mathcal{T}W)(s)|
\le \max_a \gamma \sum_{s'} P(s' \mid s,a) |V(s') - W(s')|.
\]
Using $\sum_{s'} P(s' \mid s,a)=1$,
\[
\le \gamma \|V - W\|_\infty.
\]
Taking supremum over $s$ yields the result.
\end{proof}

\begin{corollary}
The Bellman operator has a unique fixed point $V^*$, and iterative updates
\[
V_{k+1} = \mathcal{T}V_k
\]
converge to $V^*$.
\end{corollary}

\subsection{Worked Example: A Finite MDP}

Consider an MDP with two states $\{s_1, s_2\}$ and one action per state:

\begin{itemize}
\item $R(s_1) = 1$, $R(s_2) = 0$
\item $P(s_2 \mid s_1) = 1$
\item $P(s_2 \mid s_2) = 1$
\end{itemize}

Then:
\[
V(s_2) = 0 + \gamma V(s_2) \Rightarrow V(s_2)=0,
\]
\[
V(s_1) = 1 + \gamma V(s_2) = 1.
\]

\medskip

\noindent
\textbf{Interpretation.}  
The value reflects immediate reward plus discounted future reward.

\subsection{State-Transition View}

An MDP can be visualized as a directed graph:
\begin{itemize}
\item nodes: states
\item edges: actions and transitions
\item edge weights: rewards and probabilities
\end{itemize}

\medskip

\noindent
\textbf{Interpretation.}  
Reinforcement learning is path optimization over this graph under uncertainty.

\subsection{Connection to Representation Learning}

In high-dimensional settings, raw observations $x$ may not satisfy the Markov property.  
A representation $\phi(x)$ is often learned such that
\[
P(s_{t+1} \mid s_t, a_t)
\approx
P(s_{t+1} \mid x_t, a_t).
\]

\medskip

\noindent
\textbf{Insight.}  
Representation learning in RL aims to construct state variables that make the environment approximately Markov.

\subsection{What to Remember}

\begin{itemize}
\item RL models sequential decision making under uncertainty
\item Value functions satisfy fixed-point equations
\item The Bellman operator is a contraction
\item Representations in RL aim to recover Markov structure
\end{itemize}

\subsection{Common Confusion}

\begin{itemize}
\item The Markov property is a modeling assumption, not always satisfied
\item Value functions depend on policies
\item Optimality equations are nonlinear due to the max operator
\end{itemize}

\subsection{Exercises}

\begin{exercise}
Compute the value function for a three-state chain with deterministic transitions.
\end{exercise}

\begin{exercise}
Show that policy evaluation defines a linear system in $V^\pi$.
\end{exercise}

\begin{exercise}
Prove that the Bellman expectation operator is also a contraction.
\end{exercise}

\begin{exercise}
Construct an example where a poor representation violates the Markov property.
\end{exercise}

\subsection{Notes and References}

Key sources include:
\begin{itemize}
\item classical dynamic programming and optimal control
\item foundational RL texts (e.g., Sutton and Barto)
\item modern deep RL and representation learning work
\end{itemize}

This area is comparatively well understood in finite settings, but becomes substantially more complex in large-scale or partially observed environments.